% =====================================================================
%  Harald Høffding, Montaignes Betydning i Etikens Historie
%  In: Det nittende Aarhundrede. Maanedsskrift for Literatur og Kritik, udgivet af Georg Brandes og Edvard Brandes.  [Aargang 1876] April--September.  Juni 1876, pp. 210-230.
%  Antiqua.
%
%  TRANSCRIPTION -- GENERATED; DO NOT EDIT BY HAND.  Høffding's four pieces of 1875-76 (Herbert Spencers Filosofi; the
%  review of Wilkens; Montaignes Betydning i Etikens Historie; Benedikt Spinoza) are built by ONE pipeline
%  (ebook-method/hoeffding-1875-76/) into a combined working file, which split.py cuts into four transcriptions.
%  Rebuild from the repo root:   sh ebook-method/hoeffding-1875-76/rebuild.sh
%
%  STATUS: COMPLETE (2026-10-15): the volume's title page and the whole text, pp. 210-230, with its 12 notes.
%
%  EDITION.  The only printing.  Text of record: the Internet Archive's own Scribe scan of the University of Illinois copy, detnittendeaarh1876unse_0 (500 ppi colour JP2); extract ~/bibliotek/Høffding, Harald/1876-nittende-aarhundrede-apr-sep-pp210-230.pdf (SCANS.tsv, sha256), made by ia_extract.py
%  (title leaf, Indhold leaves, the article).  SECOND COPY (voter only): Google's scan of the Harvard copy, nittendeaarhund03unkngoog (its ABBYY OCR).
%  The article head and the running heads read „Etikens“; the volume's Indhold „Ethikens Historie“ (both recorded; the catalog follows the head).
%
%  ACCURACY (2026-10-15; one collation for all four texts, 57 working pages, ~17,100 words).
%  (a) The DRAFT is our own tesseract reading (tessdata_best: dan for Det nittende Aarhundrede, script/Fraktur for Nær og
%      Fjern; 300 dpi) of the Archive's JP2 images, laid over the working PDF and turned into LaTeX by draft.py.  It was
%      chosen over the Archive's own OCR (tesseract 5.2/5.3, hOCR) by measurement against the Harvard copy's ABBYY OCR
%      (words unmatched: 3.0% against 9.5% in the Montaigne volume; the Archive reads aa as åa and long s as f).  The
%      Archive's OCR is the WITNESS: collated word by word (collate/check.py): 900 disagreements, 76 settled by check.py.
%  (b) Rules (rules.py: the Archive's å for a, long s read f, glyph faults against attested words, 326) and a THIRD reading
%      voted: the ABBYY OCR of ANOTHER COPY (Google's scan of the Harvard volume), by context and by word alignment
%      (rules.py harv-*, vote2.py hv-*), never on figures and never making a word the corpus does not attest.  AUDIT: 30
%      letter-changing rule/vote decisions read blind at the image: 0 contradicted (6 undetermined crops).
%  (c) The rest of Det nittende Aarhundrede read at the image BLIND (round b1: 210 crops on 12 sheets, five readers, 23/24
%      planted controls right; ~80 decided by hand from the readers' lines).  OPEN 0.
%  (d) Every FOOTNOTE read whole at the image (readers N1, N2: 26 notes on 23 pages), set from the reading; the note marks
%      checked at zooms (four the draft had misplaced, moved).  STRUCTURE read at the image (reader K): the title pages, the
%      heads, the endings and signatures, the verse on p. 116.
%  (e) Nær og Fjern (FRAKTUR, two columns, no second copy): each column PROOFREAD WHOLE at the image against the draft by
%      two readers (P1, P2; 3 of 4 planted errors found): 171 corrections (proofapply.py); then an unattested-word sweep
%      and a long-s / ø sweep, every doubtful word read at a zoom (proof-fixes.tsv, final-fixes.tsv).
%  (f) LETTERSPACING (\emph{} here) MEASURED at the image (spacemeasure.py, white runs between letters, threshold 5 px: 31
%      runs), every run judged at the image by reader K (13 kept in the body; the notes' spacing read by N1/N2).  Italics:
%      none found by stroke slant (check.py --italics).
%  (g) QUOTATION MARKS: the print quotes „...“ only (every quotation read at the image in b1, the notes and the structure);
%      both OCRs read „ as , or ,, and “ as ' ” " »: set by class (quotes.py), outside the notes.
%  (h) Sweeps on the collated text, every doubtful spot read at a zoom: unattested words (229 words in no other
%      transcription: 12 words mended, the misprints below confirmed); specks both OCRs had read as marks (specksweep.py,
%      marks the Harvard copy lacks mid-sentence: 16 removed).  The TRANSLATORS' odd spots (64) were read at zooms before
%      the translation was finished: 22 more faults mended (specks, ɔ: read as 2:, a running head and a footnote left in
%      the text, two paragraph breaks at page turns), 3 misprints confirmed (final-fixes.tsv).
%  (i) SAMPLED PAGE CHECK (2026-10-15, budget-limited to 3 pages): pp. 113 (Spencer), 217 (Montaigne) and one Nær og Fjern
%      column (p. 2, right), chosen at random (seed 20261015, one per stratum), transcribed BLIND at the image by one
%      reader (1,143 words) and compared word by word (pagecheck.py): 10 differences; decided at zooms: 3 reader slips;
%      7 errors of the collated text (2 words: „for“ for „før“, „skikkelig“ for the print's „stikkelig“; 5 marks: four
%      specks both OCRs had read as stops or quotes, one stop hidden by a planted control) = ~6 per 1000 words, all mended.
%      The specks led to specksweep.py over the whole of Det nittende Aarhundrede: 16 more found at zooms and removed.
%  PRINTED AS IS (misprints kept, marked at the line):
%    p. 227: „kaldter“ for kaldte; sweep zoom
%    p. 219: „egen Væsen“ for eget Væsen; translator odd spot, zoom
%    p. 223: „kan ikke formulerer“ kan for han; translator odd spot, zoom
%    p. 227: „uudertrykt“ for undertrykt; note reader N2, translator odd spot, zoom
%  CONVENTIONS.  Page turns: % --- p. N --- and \opage{N} (the printed page); a word divided over a page turn is
%  split with %.  Footnotes as \footnote at their marks, numbered per page as printed (*), **)).  \emph{} = letterspacing.
%  The running heads and the sheet signatures are not transcribed.  The ø rule (CLAUDE.md).
% =====================================================================
\documentclass[12pt,a4paper,openany]{book}
\usepackage[utf8]{inputenc}
\usepackage[T1]{fontenc}
\usepackage{libertinus}
\usepackage{libertinust1math}
\usepackage[danish]{babel}
\usepackage[protrusion=true,expansion=true]{microtype}
\usepackage{setspace}
\usepackage[margin=1.3in]{geometry}
\usepackage{fancyhdr}
\usepackage{hyperref}
\hypersetup{pdftitle={Montaignes Betydning i Etikens Historie}, pdfauthor={Harald Høffding}, hidelinks}
\pagestyle{fancy}
\fancyhf{}
\fancyhead[LE,RO]{\thepage}
\fancyhead[C]{\textit{H. Høffding: Montaignes Betydning i Etikens Historie (1876)}}
\fancypagestyle{plain}{\fancyhf{}\fancyhead[LE,RO]{\thepage}}
\setstretch{1.3}
\usepackage{marginnote}
\usepackage{graphicx}
\DeclareUnicodeCharacter{0254}{\reflectbox{c}}   % ɔ (id est)
\makeatletter
\renewcommand{\thefootnote}{\ifcase\value{footnote}\or *\or **\or ***\else\arabic{footnote}\fi)}   % the print marks each page's notes *), **)
\makeatother
\newcommand{\opage}[1]{\marginnote{\footnotesize\textit{[#1]}}}
\newcommand{\signature}[1]{\par\smallskip\hfill #1\hspace*{2em}\par}
\pdfstringdefDisableCommands{\renewcommand{\footnote}[2][]{}\renewcommand{\opage}[1]{}}

\begin{document}

\begin{titlepage}
\begin{center}
{\LARGE\bfseries Det nittende Aarhundrede.}\\[1em]
{\small\bfseries \emph{Maanedsskrift}}\\[0.3em]
{\small for}\\[0.3em]
{\large\sffamily Literatur og Kritik.}\\[1.2em]
{\small Udgivet af}\\[0.3em]
{\large Georg Brandes {\small og} Edvard Brandes.}\\[0.6em]
\rule{3em}{0.4pt}\\[0.6em]
{\small\bfseries 1876. April---September.}\\
\vfill
% [printer's device: a ring with the monogram GB]
{\large\sffamily \emph{København}.}\\[0.4em]
{\small Forlagt af den Gyldendalske Boghandel (F. Hegel).}\\
{\small Trykt hos J. Jørgensen \& Co.}\\[0.4em]
{\small 1876.}
\end{center}
\end{titlepage}
\clearpage

% --- p. 210 ---
\setcounter{footnote}{0}%
\opage{210}\begin{center}{\Large\sffamily\bfseries Montaignes Betydning i Etikens Historie.}\\[0.5em]\rule{4em}{0.4pt}\end{center}\medskip

Der kan drages en Parallel mellem Videnskabens
Genfødelse i det 15de og 16de Aarhundrede og det
Vendepunkt, som indenfor den græske Tænknings Historie
indtraadte med Sokrates. De græske Filosofer før ham
havde fordybet sig i Betragtningen af Verden, den ydre
Natur, som Grækerne kaldte Himlen. Sokrates indskærpede
nu, at for at Mennesket kan have en Viden,
som kan hjælpe ham i Livet og blive Sjælen i hans
Handlen, maa han først og fremmest kende sig selv.
Det Menneskelige, Menneskets Natur, Opgaver og Pligter
var for Sokrates Tænkningens væsenlige Genstand, og
han er derved bleven Grundlægger af Etiken, Læren
om den menneskelige Handlens Vilkaar og Formaal. I
Forhold til de tidligere Tænkere stod han da som den,
der førte Filosofien ned fra Himlen til Menneskenes Huse
og Boliger. Ogsaa om Renæssancetidens Tænkere kan
dette Billede bruges. Men Himlen betegner da ikke,
som hos Grækerne, den ydre Natur, Universet, men den
overnaturlige Verden, i hvilken Middelalderens Kirke
havde lært Menneskene at se Livets egentlige Maal, i
Forhold til hvilket det jordiske Liv selv var et Intet.
Verdenslivet havde ingen selvstændig Betydning, ingen
væsenlige Opgaver at løse; det var en kort Prøvetid,
som afsluttedes med en absolut Afgørelse. Kunstens og
Videnskabers Genfødelse var Tegn paa, at man atter
% --- p. 211 ---
\setcounter{footnote}{0}%
\opage{211}begyndte at vende tilbage fra Himlen til Menneskeverdenen.
Man hentede især Styrke ved Betragtningen
og Tilegnelsen af den antike Kultur, hvor man fandt det
rent menneskelige Liv afspejlet, det „naturlige Menneske“,
som Kirken fordømte. Naturen i det Hele blev Genstand
for begejstret Anskuelse, og den fremblomstrende Naturvidenskab
gav snart denne Anskuelse et fast og urokkeligt
Grundlag. Man følte sig ikke længere fremmed i
Universet; der opstod en Frejdighed og en Tillid til
Livet, som gjorde en Stræben efter Løsning af rent
menneskelige Opgaver mulig. Men denne Tillid viste sig,
som Følge af Tidsalderens hele Karakter, paa umiddelbar,
tumultuarisk Vis, som brændende Entusiasme, ikke
med klar Selvbesindelse. Man ser endnu Menneskets
Væsen som indflettet i Natursammenhængen, opfatter det
ikke i hele dets Dybde og Inderlighed. Men Naturen
staar som den uendelige, alt omfattende, hvori alle endelige
Modsætninger tabe sig, navnlig Oldtidens Modsætning
mellem Form og Stof, Middelalderens Modsætning mellem
Aand og Natur.

Hvad særlig Etiken angaar, da fremdrog den nu
paa den ene Side det, som Kirken kaldte det naturlige
Menneske, med dets Drifter og Stemninger, dets Tanker
og Kræfter, som det virkelige Grundlag for Nydelse og
Handlen (Montaigne, Macchiavelli) eller for at vise, at
det i sig selv har Principet for den sædelige Handlen
og ikke behøver at vente Lønnen i et Hinsides (Pomponatius);
paa den anden Side finder man Frigørelsen i
Hengivelse til og begejstret Opofrelse for det uendelige
Ideindhold, som nu opgaar for Bevidstheden (Giordano
Bruno). Naturligvis kunne disse to Sider ikke adskilles;
de staa i indre Forbindelse med hinanden; det følger
allerede af de foregaaende Bemærkninger, og det vil
navnlig indlyse ved den følgende Skildring af En af de
nævnte Tænkere fra Renæssancetiden.

% --- p. 212 ---
\setcounter{footnote}{0}%
\opage{212}Montaigne er en typisk Repræsentant for den Skepticisme,
der i Tider fulde af stærke aandelige Modsætninger
ofte opstaar hos klare, uhildede Aander, der ikke
føle Trang til en positiv Anskuelse, allerede af den
Grund, at man jo maatte træde i Marken for en saadan
og altsaa selv blive en af de mange kæmpende Parter.
Det vilde Montaigne for enhver Pris undgaa. Han levede
i Frankrig (1533---1592) paa Religionskrigenes
Tid og var her Vidne til, hvorledes enhver af de stridende
Parter ikke blot mente at have Retten paa sin
Side, men var sikker paa at besidde den eneste saliggørende
Sandhed, saa sikker, at man med Ild og Sværd
vilde bringe Modstanderne over paa sin Side. Montaignes
sunde og selvstændige Natur --- man har kaldet ham
det 16de Aarhundredes mest originale Aand --- udviklet
ved en forstandig og naturlig Opdragelse, ved Rejser og
ved omfattende Læsning især af klassiske Forfattere,
maatte nødvendigvis føle sig foranlediget til at undersøge
den menneskelige Erkendelses og Overbevisnings
Væsen, der kunde tillade en saadan Forskellighed og
Strid. Dog maa man her ikke vente at finde metodiske
Undersøgelser. Tænkningen er for ham snarere som et
Slags Leg eller Tidsfordriv; han lader Tankerne gaa,
som de ville, og giver atter sine Læsere dem ligesaa
ubundne og uordnede, som han selv har faaet dem\footnote[1]{„Je n'ay pas aultre sergeant de bande à renger mes pieces que la fortune: à mesme que mes resveries se presentent, je les entasse“. Essais de Michel de Montaigne liv. II chap. 10.}.
Derfor lover han heller ikke sine Læsere Lærdom; det
Eneste, de ville lære at kende af hans Bog, er ham selv:
C'est icy un livre de bonne foy, lecteur. Il t'advertit dez
l'entré, que je ne m'y suis proposé aulcune fin, que domestique
et privée“. (L'aucteur au lecteur),

Men allerede dette er betydningsfuldt. Saaledes at
hengive sig til sine egne Tanker og vugge sig i sine
% --- p. 213 ---
\setcounter{footnote}{0}%
\opage{213}egne Fantasier er et moderne Træk, et Vidnesbyrd om,
at Individet allerede har frigjort sig fra de traditionelle
Forudsætninger og uden at bekymre sig om nogen Autoritet
trygt tør overlade sig til Strømningerne i sit eget
Indre og lade sig lede af dem paa Opdagelsesrejser til
nye, ubekendte Egne. Allerede dette giver os derfor
Anledning til at spørge om Montaignes Forhold til
Kirkelæren.

I sin Ungdom havde Montaigne efter sin Faders
Opfordring oversat et Skrift af Raimundus af Sabunde
(som levede i det 15de Aarhundrede), en naturlig Theologi,
som i Naturen og Fornuften søgte Beviser for de
kirkelige Lærdomme. Senere hørte han, at der fra forskellige
Sider var rejst Indvendinger mod denne Bevisførelse,
og i det længste Kapitel blandt hans „Essays“
(liv. II chap. 12) søgte han saa at imødegaa disse Indvendinger.
Men han gør dette paa en Maade, der trods
den gennemførte Ironi røber hans virkelige Anskuelse.

Paa den ene Side havde man ment, at Troeslærdommene
vare for hellige og ophøjede til at støttes med
menneskelige Grunde; paa den anden Side havde man
erklæret de anvendte Grunde for svage og ubrugelige.

De første Modstandere viser han sig tilsyneladende
meget velvillig imod. Meget vel, siger han. Troen skal
i sig selv være nok. Men --- hvis vor Tro virkelig var
ægte, indgiven fra Himlen, saa vilde vort Liv sikkert se
anderledes ud; alle Forhold vilde være omdannede i
Overensstemmelse med de hellige Fordringer, og de
Kristne vilde ikke saa ofte blive beskæmmede i deres
Vandel af Tyrker og Hedninger. Saa vilde heller ikke
Religionskrigene med alle deres uædle Lidenskaber have
kunnet splitte os; den menneskelige Affekt og Interesse
vilde ikke have ført Ordet ved disse hellige Anliggender.
Alt dette er Bevis nok for, at vi stedse tage Religionen
paa vor egen Maade, med vore menneskelige Hænder, og
at den overnaturlige Indgivelse træder tilbage for Tradition,
% --- p. 214 ---
\setcounter{footnote}{0}%
\opage{214}Skik og Brug. Vi ere Kristne af lignende Grunde som
vi ere Franske eller Tyske.

Allerede her har Montaigne glemt Raimundus, som
han skulde forsvare. Men han taber ham aldeles af
Syne, saasnart han giver sig i Kast med den anden
Indvending, den, at Beviserne ikke slaa til. Disse Modstandere
vil han aldeles knuse ved at vise dem den
menneskelige Erkendelses Afmagt og Intethed. Men
istedetfor den Slutning, som man skulde vente heraf, at
Mennesket blindt maa hengive sig til den guddommelige
Avtoritet, skinner det klart igennem, at Montaigne selv
akcepterer de skeptiske Resultater, og at de for ham
danne en Grænse ikke blot overfor den filosofiske, men
ogsaa overfor den teologiske Dogmatisme. Han hæver
det Guddommelige saa højt over Menneskets Blik og
Evne\footnote[1]{C'est aux chrestiens une occasion de croire, que de rencontrer chose incroyable. Liv. II. chap 10.}, ikke for at tilbede det des inderligere, men for
med Ironi at vende sig bort derfra.

Montaigne stiller de forskellige filosofiske Skolers
Lærdomme op imod hverandre; han viser Sansernes Upaalidelighed;
han samler Træk af Dyrenes Liv for at godtgøre,
at Mennesket ikke besidder noget væsenligt Fortrin
fremfor dem; --- kort sagt, han har hele Skepticismens
sædvanlige Arsenal til sin Raadighed, og han
anvender det med Virtuositet. Men hvad der især maa
interessere os, er at se hvorledes han forstaar at koncentrere
de skeptiske Argumenter paa en Maade, der
røber en virkelig filosofisk Tænker bag den kaade Ironiker.
Vor Fornuft, siger han, kan ikke vise os Sandheden;
thi for at vi kunde være sikre paa, at det virkelig
var Sandheden, den viste os, maatte vi have en anden
Fornuft, der kunde indestaa for den førstes Sandhed, og
for denne anden ogsaa en tredie og saaledes i det
Uendelige. Al vor Erkendelse stammer tilsidst fra
% --- p. 215 ---
\setcounter{footnote}{0}%
\opage{215}Sanserne\footnote[1]{Toute cognoissance s'achemine en nous par les sens; se sont nos maistres (ibid).}, men heller ikke de kunne lære os det Sande.
Thi de opfatte kun Tingen saaledes som den ifølge deres
Beskaffenhed tager sig ud for dem\footnote[2]{Les sens ne comprennent pas le subject estrangier, ains seulement leurs propres passions (ibid).}. Tingen selv undgaar
os bestandig, og vi have ingen Maalestok for at
bedømme Forholdet mellem dens virkelige Væsen og vor
Fornemmelse af den. Paa lignende Maade viser Montaigne,
at da Indbildningskraften ikke kan gaa ud over
vort eget Væsen, tænke vi os nødvendigvis Guddommen
under menneskelige Former. Skønt Menneskene altsaa
selv have dannet sig deres Guder, blive de dog bange
for dem, ligesom Børnene blive bange for deres Kammerats
Ansigt, som de selv have malet eller sværtet. ---
Der ligger her til Grund en Anelse om den Tanke, Kant
og Feuerbach i nyere Tid saa energisk have gennemført
paa Erkendelseslærens og Religionsfilosofiens Omraade,
at det Sete stedse betinges ved Synet og ikke kendes
uden i sit Forhold til dette.

Montaigne drager heraf saavel teoretiske som navnlig
praktiske Konsekvenser. Da Mening staar mod Mening
hos de Lærde, --- da der ikke er noget saa Urimeligt,
at jo et eller andet Folk har optaget det i sin Religion,
--- da det Gode og Rette er forskelligt hos de forskellige
Nationer, saa at det, der er Synd paa denne Side
Bjærgene, er Dyd paa den anden Side, --- saa er der
ikke Andet for end at vedgaa sin Uvidenhed. Det er
et Bevis paa Galskab og paa den allerhøjeste Grad af
Uvidenhed at anse sig for sikker paa Noget. Al vor
Forsken kan kun tjene til at udvikle og bekræfte den
Uvidenhed, som er os medfødt. Ja, man kan ikke engang
være vis paa sin Tvivl og Uvidenhed; man skal
derfor ikke sige: jeg tvivler! --- men: hvad véd jeg? ---
% --- p. 216 ---
\setcounter{footnote}{0}%
\opage{216}Heraf følger, at man vel kan undskylde dem, der fejle i
deres religiøse Forestillinger; thi da disse Forestillinger
angaa Noget, som ligger ud over menneskelige Grundes
Raaderum, saa kan den, der ikke oplyses af den guddommelige
Naade, let fare vild; --- men at det paa den
anden Side er uforsvarligt at sætte sig op imod den bestaaende
Tilstand, fordi man mener selv at have fundet
Noget, som er bedre. Det vidner om stor Egenkærlighed
og Hildethed. Jeg væmmes ved det Nye, siger Montaigne,
under hvilken Skikkelse det saa fremstiller sig. Ikke
fordi de bestaaende Love i sig selv altid ere fornuftige.
Lovenes Gyldighed beroer ikke paa, at de ere retfærdige,
men netop paa, at de ere Love: c'est le fondement mystique
de leur auctorité! Ofte skriver Lovenes Oprindelse
sig fra Taabelighed og Egennytte, men Sædvanen
giver dem Fasthed og lader dem synes retfærdige. Derfor
er det farligt at følge dem tilbage til deres Oprindelse.
Det gaar med dem, som med en mægtig Flod,
der i sit Udspring kun er et ubetydeligt Væld. Saaledes
vokse Lovenes Magt fra en tilfældig Begyndelse i
Tidens Løb.

Sædvanen er Verdens Herskerinde. Det gælder ikke
blot om de juridiske Love, men ogsaa om de moralske.
Samvittigheden føler sig bunden til, hvad den er vant
til at respektere, og kan ikke uden Anger og Nag frigøre
sig derfor. Kannibalerne, der betragte det som en
Pietetspligt at spise deres gamle Forældre, vilde forarges
ved at se de Afdøde blive nedgravede i Jorden til
Føde for Ormene. Vel skal den Vise indadtil frigøre
sin Sjæl, saameget som muligt drage sig ud af Mængden
for uhildet at kunne dømme om Tingene (retirer son
ame de la presse et la tenir en liberté et puissance de
juger librement des choses. Liv. I chap. 22); --- men udadtil
maa han følge de bestaaende Love og Skikke. I
Modsætning til de mange stridende Meninger overlader
han sig umiddelbart og naivt til Naturen, som aldrig
% --- p. 217 ---
\setcounter{footnote}{0}%
\opage{217}leder fejl, naar man blot ikke forfalsker den ved at ville
tolke dens Visdom; men herimod sikrer netop Uvidenheden
og Ligegyldigheden ham: „Oh! que c'est un doulx
et mol chevet, et sain, que l'ignorance et l'incuriosité, à
reposer une teste bien faicte!“ (liv. III chap. 13). Fra
de udvortes Sysler trækker han sig tilbage i sit eget Jeg
og følger det ind i dets inderste Løngange, de mørkeste
Folder. Men denne Vej til Ro kan saa naturligvis ikke
være fælles for Alle; Enhver maa søge den for sig, efter
sin Natur.

Denne gennemførte Skepsis, der ikke engang i sin
egen Tanke finder en fast Grund, men bestandig holder
sig i Spørgsmaalsform, er dog ikke Montaignes sidste
Ord. Allerede Pascal har set dette, naar han i den
mesterlige Karakteristik af Montaigne, som findes i hans
„Pensées“, sætter det som dennes Formaal at finde en
Moral, der grundes paa Fornuften, uden Troens Lys.
Montaigne er en Moralist, ikke blot fordi han psykologisk
 skildrer Menneskets Natur, men ogsaa fordi han
virkelig udtaler en sammenhængende etisk Anskuelse.
Men af hvad Art kan denne være, naar han stiller sig
skeptisk overfor ethvert positivt Princip?

Vi have allerede naaet det Punkt, hvor Skepticismen
hos Montaigne gaar over til positiv Anskuelse. Han
hævder vel Uvidenheden og gør Que sais-je? til sit
Princip; men Uvidenheden er ham kun Middel til at
gribe Naturen. Menneskene have, siger han, ved deres
Ræsonnementer forkvaklet den sande og oprindelige Natur,
hvis Spor kun findes hos de Udannede og hos Dyrene.
Naturen har derved mistet sin faste og universelle
Karakter og har faaet en særegen Form hos hver Enkelt.
Dyden bestaar ikke i Tankens eller Fornuftens Herredømme,
men i at lade Naturen udvikle sig efter sin egen
Lov, saaledes at Fornuften blot afskærer de falske Udvækster.
Drifterne følge deres eget naturlige Løb, og
Fornuften kan ikke umiddelbart hindre dem deri. Mon%
% --- p. 218 ---
\setcounter{footnote}{0}%
\opage{218}taigne siger derfor, at hvis han besidder nogen Dyd, saa
er det en Uskyldighed, han snarere skylder Lykken end
Fornuften. Naar Naturen af sig selv fører den etiske
Harmoni med sig, saa følger det af sig selv, at Nydelsen
ikke staar i Modsætning til Dyden, men netop ved Dyden
opnaar sin højeste og reneste Form. Af Naturen tragte
vi ikke efter Andet end Nydelse. Men den lavere Nydelse
fortjener ikke dette Navn; den er kun øjeblikkelig
og fuld af Ængstelse og Savn. Den højere Nydelse bestaar
i selve Dyden, som man derfor ikke med Stoikerne
maa tænke sig mørk og stræng, tronende paa en stejl
Klippe i en ensom Ørken, der er fuld af Farer. Ej
heller forholder det sig saaledes, at den selv er mild og
smilende, men Vejen til den mørk og besværlig; saa
vilde Ingen besidde den, da Ingen naar at tilbagelægge
hele Vejen. Vejen maa svare til Maalet, hvortil den
fører. Dyden boer paa en skøn og blomstrende Høj,
hvorfra den ser alle Ting under sig, og man naaer dertil
ad skyggefulde og blomstersmykkede Stier, med Naturen
som Fører, Lykken og Nydelsen som Ledsagere.

Denne Lære minder om Epikuræismen. Men allerede
Montaigne selv har bemærket, at det kun er en falsk
Mening, at Epikuræismen i Fasthed og Strænghed skulde
staa tilbage for Stoicismen (liv. II chap. 11). Hvis man
anvender dette paa hans egen Lære, vil det kunne udtrykkes
saaledes, at der kræves meget Arbejde og megen
Anstrængelse, for at den sande Natur kan udfolde sig i
Mennesket. Hvad der gør, at Stoikerne staa Montaigne
fjærnere, er det asketiske Element hos dem og deres
Ringeagt for det umiddelbare Naturlige og Menneskelige,
som Montaigne netop søger at fremlokke og frede om.
Aandens Storhed, siger han, bestaar ikke saameget i at
hige opad og fremad, som i at ordne og begrænse sig.
Intet er saa skønt og ret, som at være Menneske paa
god og tilbørlig Vis, ingen Videnskab saa vanskelig, som
at forstaa at føre sit Liv paa rette Maade, og den værste
% --- p. 219 ---
\setcounter{footnote}{0}%
\opage{219}Sygdom er Foragt for vort eget Væsen\footnote[1]{Entre nous --- siger han i Anledning af de kristelige Asketer --- ce sont choses que j'ay tousjours veues de singulier accord, les opinions supercelestes, et les moeurs soubterraines. Liv. III chap. 13.}. Naturen er
en god Fører, og ikke mere god end klog og retfærdig.
Derfor beundrer Montaigne alligevel „de heroiske Sjæles
Højhed“. Det er jo hans Princip, at Enhver skal maales
med sin egen Alen, efter sin egen Natur (sur son propre
modele). Hans egen Svaghed gør ham ikke blind for
Andres Styrke; skønt han selv kryber paa Jorden, kan
han dog se Andres høje Flugt, og han anser det ikke
for nogen ringe Sag at bevare den moralske Dømmekraft
sund i en Tidsalder som hans, hvor Dyden ikke
synes at være Andet end „un jargon de college“, og
hvor de fleste stræbe at nedsætte enhver stor og ædel
Handling ved paa en subtil Maade at henføre den til
slette Motiver for at kunne faa den draget ned i deres
egen Sfære. Hine „heroiske Sjæle“ ere Mønstre, som
Naturen vælger for at vise, hvilken Højde den menneskelige
Dyd og Kraft kan naa (liv. I chap. 36).

I den Sædelighed, til hvilken Naturen fører os, naaer
 altsaa Nydelsen sin ædleste Form ved at blive Et med
Dyden. Den forsoner tillige en anden Modsætning,
mellem Egoismen og Hengivelsen til Andre. Ved at indlade
sig paa de offenlige Anliggender og give sig af med
Andres Sager, mister man let Noget af sit egen Væsen. De % PRINTED AS IS (egen Væsen: for eget Væsen; translator odd spot, zoom)
Fleste give sig umiddelbart hen, og tilhøre derfor ikke
sig selv. Selv Vellysten skal man ikke helt hengive sig
til; den gemmer Smerte paa Bunden. Man skal ikke
fordybe sig i Verden, men kun berøre den, glide hen
over dens Overflade. Det Største i Verden er at forstaa
at tilhøre sig selv. Man skal ikke binde sig til Andet
end til sin egen Natur (n'espouser que soi-mesme). Deraf
følger ikke, at man skal være Egoist og ikke bekymre
% --- p. 220 ---
\setcounter{footnote}{0}%
\opage{220}sig om Andre; tværtimod, den som ikke lever for Andre,
lever heller ikke for sig selv. Selvhensynet indeslutter
Hensynet til Andre; begge følge af den samme Natur.

Montaigne vil ikke foredrage en slap og blødagtig
Nydelseslære. Det ses ikke blot af det Foregaaende,
men ogsaa af den Rolle, Tanken om Døden spiller hos
ham. Ingen, siger han, vilde kunne glæde sig ved Livet,
naar han frygtede Døden, lige saa lidt som den dømte
Forbryder paa Vejen til Retterstedet kan glæde sig over
de skønne Egne, han kommer igennem. De Fleste søge
derfor den Udvej at se at glemme den, men denne Udvej
er dyrisk, er brutal Stupiditet. Og desuden, hvad
hjælper det? Man gaar og kommer, danser og ler, som
om Døden ikke var til; og naar den saa berøver os En
af vore Nærmeste, forandres denne Sorgløshed til den
dybeste Fortvivlelse og Nedslagenhed. I Sandhed, selv
om denne dyriske Tankeløshed var mulig hos et forstandigt
Menneske, saa er den dog for dyrt købt. Man maa
i rette Tid lære at bekæmpe denne uundgaaelige Fjende,
uforfærdet gøre Front imod den. Dens frygteligste Vaaben
er det Fremmede og pludselig Truende. Lad os da berøve
den dette ved at gøre os fortrolig med den, vænne
os til den ved stedse at have den i Tanke, ved at fremstille
den for vor Indbildningskraft under alle mulige
Skikkelser. Lad os gøre som de gamle Ægyptere, der
ved deres muntreste Fester lod bringe Mumier af afdøde
Mennesker ind for at tjene til Paamindelse for de Indbudte.
Da det er uvist, hvor Døden venter os, saa lad
os vente den overalt. At tænke over Døden er at tænke
over Friheden; den som har lært at dø, har lært at ophøre
at trælle; der gives intet Ondt i Livet for den, der
har indset, at Berøvelsen af Livet ikke er et Onde. Man
skal ikke knytte sig saa fast til Livet, ej heller til nogen
Nydelse eller Gerning i Livet, at man ikke uden Sorg
kan skilles derfra. Livet er i sig selv hverken et Gode
eller et Onde; det kommer an paa, hvorledes man har
% --- p. 221 ---
\setcounter{footnote}{0}%
\opage{221}brugt det; man kan have levet meget i kort Tid. „Je
veux, qu'on agisse et qu'on alonge les offices de la vie,
tant qu'on peult, et que la mort me trouve plantant
mes choulx, mais nonchalant d'elle et encore plus de
mon jardin imparfaict“. Hvor er desuden vort Liv kun
et forsvindende Intet, naar vi sammenligne det med Evigheden,
eller endog blot med Stjærnernes, Bjærgenes.
Flodernes og Træernes Ælde! Hvilken Ret have vi til
at tillægge os en Væren i dette korte Øjeblik, der kun
er et lille Glimt i en evig Nats uendelige Løb, en saa
kort Afbrydelse af vor vedvarende og naturlige Tilstand,
da jo Døden indtager Alt, hvad der ligger forud og bagefter
dette Øjeblik, ja endog en god Del af det selv?
Med Rette siger derfor Naturen til os: „Gaa ud af denne
Verden, som du er kommen ind i den. Den Overgang,
du foretog fra Død til Liv, uden Lidenskab og uden
Rædsel, skal du nu gøre om fra Liv til Død. Din Død
er et Led i Universets Orden, i Verdens Liv“\footnote[1]{Liv. I chap. 19: Que philosopher c'est apprendre à mourir.}.

Den Karakter, Mennesket udvikler i Løbet af sit
Liv, udspringer med Nødvendighed af hans oprindelige
Natur. Den bevidste Fornuft kan, som vi allerede have
hørt, ifølge Montaigne kun virke indirekte, afværgende
og rensende; den positive Livsfylde oprinder af Naturens
Skød. Dog er der én positiv Følelse, Fornuften kan bevirke,
og som opstaar, naar Virkningen af den uendelig
omskiftende Naturdrift kommer i Strid med den fornuftige
Indsigt i det Rette. Denne indre Disharmoni
føles som Anger. Herved oprettes der en indre Domstol
i os, som Ingen kan undfly, selv om han overfor
Andre synes nok saa rolig og sikker med Hensyn til
sine Handlinger. Selv om vore Gerninger tage sig glimrende
ud for Andre, ere de dog for os selv kun som
klare Straaler, der springe i Vejret fra en uren og mudret
Grund. Men skønt Angeren er en positiv Følelse, har
% --- p. 222 ---
\setcounter{footnote}{0}%
\opage{222}den dog sin Grænse der, hvor Fornuften har sin; selve
den Lidenskab, der er Et med vort eget Væsen, og som
ved Vane og Øvelse har fæstet endnu dybere Rod, kan
den ikke forandre. I denne Forstand er Angeren umulig.
Det er kun Illusion eller Hykleri, naar man mener virkelig
at kunne blive eller være bleven en Anden. Det
sker maaske i det Højeste, som Pythagoras mente, at
Menneskene faa en ny Sjæl, naar de nærme sig Guderne
for at modtage deres Orakler; men denne nye Sjæl er
kun paatagen og forsvinder igen efter kort Tids Forløb.
Enhver har i sig selv en ham ejendommelig, hans Væsen
beherskende Form (une forme sienne, une forme maitresse),
som ikke kan forandres, men modstaar enhver
ydre Paavirkning og omdanner den paa sin Vis. Derfor
naa Forsøgene paa at reformere og forbedre Verden kun
til Skallerne, aldrig til selve Kærnen. Hvad mig selv
angaar, siger Montaigne, da kan jeg ønske, at jeg var
et andet Menneske; jeg kan fordømme og afsky min almindelige
Natur (ma forme universelle) og bede Gud
om en fuldstændig Reformation af mit Væsen; men det
kalder jeg ikke Anger, ligesaa lidt som jeg kalder det
Anger at være utilfreds over ikke at være en Engel eller
en Cato. Fordi jeg forestiller mig højere og fuldkomnere
Naturer end min egen, derfor forbedres ikke mine
Evner,  ligesaa lidt som min Arm eller min Forstand
bliver kraftigere, fordi jeg kan danne mig en Forestilling
om en Anden, hos hvem de ere kraftigere. I saa Fald
vilde jo ogsaa vore mest uskyldige Handlinger blive
Genstand for Anger, fordi vi let indse, at de under Forudsætning
af en fuldkommen Natur ogsaa vilde være
blevne udførte med større Fuldkommenhed og Værdighed\footnote[1]{Liv. III chap. 2: Du repentir.}.

Altsaa: den virkelige Anger, Angeren i strængeste
Betydning vilde føre til, at jeg blev en Anden; men
% --- p. 223 ---
\setcounter{footnote}{0}%
\opage{223}en Anden kan jeg ikke blive, saalænge jeg er mig
selv. Det er denne Modsigelse i Angerens Begreb,
hvorpaa Montaigne støtter sin Benægtelse af dens
Realitet. Det ses her klart, hvilken Rolle Immanentsprincipet
spiller her hos ham, skønt kan ikke formulerer % PRINTED AS IS (kan ikke formulerer: kan for han; translator odd spot, zoom)
det filosofisk. I Modsætning til den ortodokse
Anskuelse om Menneskenaturens Fordærvelse og
Nødvendigheden af en absolut Genfødelse, hævder han
den menneskelige Naturs Enhed og Uforanderlighed og
dens Sammenhæng med Naturen i dens Helhed. Naar
Handlingerne ere udførte, gaa de ind som nødvendige
Led i Universets store Proces, i Aarsagernes og Virkningernes
uendelige Række; og det er forgæves at
sørge over dem. Fantasien kan ikke med al sin Styrke
gøre et eneste Punkt indenfor denne store Proces til Intet.
--- En saadan Anskuelse vilde, hvad f. Eks. Schopenhauers
Lære viser, kunne føre til Pessimisme; naar man
 nemlig antager, at Grundvæsenet (la forme maistresse)
er ondt. Men Montaigne betænker sig intet Øjeblik paa,
om han skal give Heraklit eller Demokrit Ret: jeg
tror ikke, siger han, at der er saa megen Ulykke blandt
os som der er Forfængelighed, eller saa megen Ondskab;
som der er Dumhed!

Pascal har skildret Montaigne som Skeptiker og
Moralist, og i det Hele som den, der fremstiller „det naturlige Menneske“. Men han har ikke ført Skildringen
videre saaledes, at det sidste Grundlag for Montaignes
Livsanskuelse kom for Dagen, det Punkt, hvorom Alting
tilsidst drejer sig, og hvorved Montaigne naaer ud over
det blot Populære og Moralske og antyder en hel filosofisk
Verdensanskuelse. Dette maatte i og for sig have
mindre Interesse for Pascal, der selv stillede sig paa et
afgjort praktisk Standpunkt. Pascal beundrer Montaignes
dialektiske Gennemførelse af Skepticismen og akcepterer
den i Paradoksets Tjeneste; men han forkaster hans
Moral som Hedenskab, ja i sidste Instans som Fejghed.
% --- p. 224 ---
\setcounter{footnote}{0}%
\opage{224}Om man nu just kan beskylde en Livsanskuelse for
Fejghed, der gør sig det til en væsenlig Opgave at indøve
Tanken om Døden og have den til sin Fortrolige
ved Alt, hvad man foretager, er en anden Sag; i sin
brændende Lidenskab for Evangeliet maatte Pascal anse
Alt for lavt, hvad der ikke absolut gik op i denne overnaturlige
Magt. Derfor er det ogsaa, at han stanser
ved den moralske og religiøse Betragtning og ikke gaar
fremad til den metafysiske. Det følger tillige af sig selv,
at det Metafysiske ligger temmelig skjult hos en causeur
som Montaigne. Men efter Alt, hvad vi have hørt i det
Foregaaende, kan det ikke være tvivlsomt, at der bag
hans Skepticisme og hans Eudæmonisme hæver sig en
uendelig Baggrund, til hvilken han stedse direkte eller
indirekte viser hen. Det er Ideen om Naturen i dens
Højhed og Uendelighed, der i sin store Kæde ogsaa omslutter
Menneskelivet som et enkelt Led. Sainte-Beuve
har i sin Port-Royal (3e éd. II p. 412), hvor han i Anledning
af Pascal kommer til at beskæftige sig med
Montaigne, peget hen paa, hvorledes denne gennem sit
Naturbegreb slutter sig til Grundtanken i de store metafysiske
Systemer, særlig Hobbes og Spinoza. Modsætningen
til Dualismen, Tilliden til Menneskenaturen og
dens Kræfter maatte nødvendigvis føre et klart og livligt
Hoved til den Grundvold, hvorpaa hele den moderne
Verdensanskuelse hviler, selv om han ikke har udtalt
denne i en filosofisk Formel.

Naturbegrebet gjorde os allerede tidligere Forbindelser
mellem Montaignes Skepticisme og hans Moral
forstaaelig. Men dette Begreb spiller saavel i det Smaa
som i det Store en Hovedrolle hos ham. Er Talen f.
Eks. om Sygdom, vil han ikke vide Noget af Lægerne;
de forkvakle kun Naturen og hindre den i dens egen
rette Udvikling; --- ligesom ethvert andet Væsen har
ogsaa Sygdommen sin bestemte Tid til at udvikle sig og
til at afslutte sit Liv, og det er forgæves at gribe ind
% --- p. 225 ---
\setcounter{footnote}{0}%
\opage{225}her. Man skal lade Naturen raade; den forstaar Sagerne
bedre end vi. I Naturens Orden maa man finde
sig. Ligesom Verdensharmonien er vort Liv sammensat
af Modsætninger, af forskellige Toner, som alle høre med
i den store Sammenhæng, Sygdom og Død saavel som
Sundhed og Liv. --- Er Talen om Opdragelse, da gælder
det ligeledes at lade Naturen have sit frie Løb: nature
peut tout et fait tout! --- Endelig hæver han ved Hjælp
af Naturbegrebet selve sin Skepticisme. Den eneste
Maade, siger han, ved hvilken man kan sikre sig mod at
maale Tingene med en falsk Alen, er at gaa ud over de
relative Sammenligninger og stille sig „vor Moder Naturen
i dens hele Majestæt“ for Øje; den viser os en
gennemgaaende og stadig Mangfoldighed, indenfor hvilken
vi opdage os selv og Alt, hvad vi kalde stort; herved
faa vi først det rette Begreb om Tingenes Størrelse.
Uden en saadan absolut Maalestok anser Enhver sin
Religion og sit Lands Forfatning for de eneste rigtige
og kalder alt Andet Barbari, medens det dog snarere
burde kaldes Barbari at indsnevre „vor store og mægtige
Moder Naturs“ Kraft ved at drage vilkaarlige og
tilfældige Grænser og Skranker. Det er desuden ogsaa
umuligt; Naturen spotter vore Kunster og triumferer over
vore kløgtige Paafund\footnote[1]{Ce n'est pas raison que l'art gaigne le poinct d'honneur sur nostre grande et puissante mere nature. Nous avons tant rechargé la beauté et la richesse par nos inventions, que nous l'avons du tout estouffee: si est ce que partout ou sa pureté reluict, elle faict une merveilleuse honte à nos vaines et frivoles entreprises. Liv. I chap. 30: Des cannibales.}. --- Og med dette Naturbegreb
staar endelig Montaignes Gudsbegreb i nøje Forbindelse.
„Af alle menneskelige Meninger om Religionen“, siger
han (liv. II chap. 12), „synes den mig at have mest Sandsynlighed
og Forsvarlighed, der erkendte Gud som en
ubegribelig Magt, Ophav til og Opholder af Alt, den
absolute Godhed og Fuldkommenhed, der tager imod og
% --- p. 226 ---
\setcounter{footnote}{0}%
\opage{226}bifalder den Ære og den Dyrkelse, som Menneskene,
hver paa sin Vis, ydede ham, under hvilken Form, under
hvilket Navn og paa hvilken Maade det saa skete.“ Vi
se heraf, hvilken nøje Sammenhæng der er mellem Montaignes
Verdensanskuelse og hans humane og tolerante
Aandsretning.  Tolerancen grunder sig hos ham ikke
blot paa Skepticismen, men Forudsætningen om, at den
evige og uendelige Natur ligger til Grund for og aabenbarer
sig i Alt og Alle, maa nødvendigvis føre til den
Erkendelse, at Ingen er afskaaret fra det Evige og Guddommelige,
og at dette kan lægge sig for Dagen under
utallige Former, som hver for sig have Ret til at udvikle
sig med Frihed.

Montaigne var ikke, som Pascal i sin jansenistiske
Iver mente, en Fejghedens og Sandselighedens Apostel,
men en ædel og sund Natur. Han afspejler sig selv i
sit Skrift baade efter sine gode og sine slette Sider, med
elskværdig Naivetet. Hvilken Inderlighed, der boede i
i ham, ses bl. A. af den Maade, hvorpaa han skildrer
sit Venskabsforhold til Etienne de la Boétie\footnote[1]{Nos ames ont charié si uniement ensemble; elles se sont considerees d'une si ardente affection, et de pareille affection, descouvertes jusques au fin fond des entrailles l'un de l'aultre, que non seulement je cognoissois la sienne comme la mienne, mais je me feusse certainement plus volontiers fié à luy de moy, qu'à moy.} (liv. I chap.
27). Dette Forhold bestod, siger han, i en fuldstændig
gensidig Opgaaen af den Enes Vilje i den Andens, saa
at Ingen beholdt Noget for sig selv, og Enhver kendte
den Andens inderste Motiver\footnote[2]{Etienne de la Boëtie (1530---1563), Medlem af Parlamentet i Bordeaux, er især bekendt ved sin „Discours sur la servitude volontaire“, ogsaa kaldet „Contre un“, som han forfattede i sit 18de Aar (par maniere d'exercitation seulement, siger Montaigne), og hvori han med glødende Begejstring, skønt i en retorisk Stil, der stærkt minder om de Gamle, forkynder Læren om Alles naturlige Frihed og viser, at det er Menneske%
\opage{227}hedens egen Skyld, naar den bliver uudertrykt, thi Undertrykkeren er kun En mod de Mange og vilde ikke have nogen Magt, naar disse ikke selv underkastede sig ham. Montaigne kalder ham Aarhundredets største Mand; senere har man ment, at det smukke Kapitel, i hvilket Montaigne omtaler sin afdøde Ven har gjort denne mere berømt, end hans egne Skrifter vilde have gjort.}. „Naar jeg sammenligner
% --- p. 227 ---
\setcounter{footnote}{0}%
\opage{227}hele Resten af mit Liv med de fire Aar, i hvilke det
har været mig givet at nyde hans kære Selskab og Omgang, saa staar det kun for mig som en Taage, som en
lang og kedsom Nat“. Venskabet var i det Hele det
etiske Forhold, der mest laa for Montaigne; her gav
han sig ganske hen uden at frygte for at miste sig selv.
--- Hvor sund en Natur han var, ser man af den Bemærkning,
at Medgang altid stemmer ham til Alvor og
Besindighed. „Jeg er frommere i Medgang end i Modgang.
Lykken tjener mig til Tugt og Opdragelse, ligesom
Ulykkerne og Revselserne for Andre; Heldet er mig
en særlig Spore til Maadehold og Beskedenhed“. (Liv.
III chap. 9).

Med Nødvendighed maatte vi vende tilbage til Personen
selv, efter tilsyneladende at have forladt ham for
at fremstille hans Livsanskuelse. Hos Montaigne mindre
end hos nogen Anden kunne begge skildres hver for sig.
Men tage vi nu hele Fænomenet, da er det en karakteristisk
Type for Renæssancen, vi her have for os.
Fuldblodig og sanselig, og dog paa sin Vis rettet mod
det Uendelige, bevægende Alt, og dog hvilende og nydende
midt under Bevægelsens Hvirvel, glædende sig ved
Oldtiden og dog med Fremtidens Tanke i sit Skød. I
sin egen Tid blev Montaigne kun forstaaet af Enkelte\footnote[1]{Se \emph{Nisard}: Précis de l'histoire de la littérature française. p. 211 f.}% PRINTED AS IS (uudertrykt: for undertrykt; note reader N2, translator odd spot, zoom)
.
Justus Lipsius kaldte ham den franske Thales, Kardinal
du Perron kaldter hans Bog le bréviaire des honnétes % PRINTED AS IS (kaldter: for kaldte; sweep zoom)
gens. I det 17de Aarhundrede stod Montaigne for Pascal
% --- p. 228 ---
\setcounter{footnote}{0}%
\opage{228}
og Jansenisterne som den store Hedning, der klarere end
nogen Anden har udtalt det naturlige ɔ: det syndige
Menneskes Principer. Først i det 18de Aarhundrede
kunde han finde beslægtede Aander. Voltaire stammer
fra ham ved sin fine og smidige Ironi og skarpe Kritik.
Rousseau udtaler, hvad Montaigne kun antyder, Nødvendigheden
af at forlade det forfinede Kulturliv (Montaignes
„nature sophistiquée et raffinée) og vende tilbage
til den ægte, oprindelige Natur. Diderot forener endelig
begge Elementer i sig og er derfor maaske den senere
Forfatter, der har størst Lighed med Montaigne.

Vi have forsøgt at finde et System i Montaignes
Essays; og med Systemet maa nødvendigvis følge en
Kritik. Som vi have søgt at vise, er Montaignes Hovedtanke
den samme, der trænger sig frem i hele Renæssancens
Tænkning, Naturens Ide. Denne Ide fremtræder
hos ham under en ejendommelig Form. Den danner den
almindelige Baggrund, den sidste Forudsætning for hans
Skepticisme; den er tillige den positive Magt, han ser
virke i Alt, og for hvilken han overalt søger at berede
Plads. Men der er dog, særlig naar vi betragte ham
som Etiker, en afgørende Mislighed i hans Opfattelse
af Naturbegrebet. Han ser det Naturlige som en umiddelbart
henflydende Strøm, som en Slags Vegeteren, der
hindres og forvanskes ved den bevidste Tænken og Virken.
Deraf bliver igen Følgen, at Menneskets Tænken og
Virken kun faar negativ eller indirekte Betydning, istedetfor
at disse Funktioner dog ligesaavel som enhver anden
maa have deres Plads og Betydning indenfor Naturen.
Principet for den menneskelige Handlen maa netop findes
ved at undersøge, hvilken Anvendelse af disse Funktioner
der følge af Menneskets sande Natur. De staa altsaa
ikke udenfor den uendelige Naturs Strøm, men ere selv
Former for denne. Gennemføres denne Betragtning, da
træder det aktive Element i Naturbegrebet klarere frem;
istedetfor den umiddelbare Strømning træder nu den
% --- p. 229 ---
\setcounter{footnote}{0}%
\opage{229}bevidste Hengivelse og den energiske Udvikling af den i
Naturideen antydede Opgave. Fra denne Side finde vi
tildels Sagen opfattet hos Giordano Bruno, Montaignes
yngre Samtidige, naar denne ligesom Platon ser Naturdriften
som Udtryk for en ideal Uro, der ikke tillader
Mennesket at hvile, førend det har erhvervet Harmoni og
Forsoning ved Tankens og Villiens Arbejde og gennem
Hengivelse til det Uendelige med den hele Personlighed\footnote[1]{Heraf følger naturligvis ogsaa en anden Opfattelse af de afgørende etiske Problemer. Saaledes fastholder Bruno navnlig Angerens Realitet, hvilken Montaigne paa naturalistisk Vis nægtede. Men Bruno opfatter Angeren saaledes, at han ikke opgiver Immanensprincipet. Han antager intet absolut Ondt, ingen Fordærvelse af Menneskenaturen. Men uden Strid og Kamp, uden Modsigelse og Modstand kan det Sande og Gode ikke sejre: og det Onde er den negative Magt, der rejser Modsigelsen og gør Kampen nødvendig. Denne negative Magt maa ophæves; Sjælen vinder sig selv gennem Sønderknuselsen. Angeren er altsaa et Tegn paa Sjælens ideale Natur: den er som en lysende Gnist, der slaar ud af den haarde Flint og vidner om dens Slægtskab med Solen. --- Det er klart, at det Principielle og Berettigede i Montaignes Polemik ikke benægtes ved Brunos Opfattelse; thi hin Polemik var rettet mod Angeren som fuldstændig Forvandling af Naturen.}.
Hos Bruno findes der derfor ingen saadan Modsætning
som hos Montaigne mellem den endelige Tænken og
Villen, som Intet kan udrette, og Anskuelsen af den
uendelige Natur. Montaignes Skepticisme, som fra en
anden Side set er et af hans Hovedfortrin, idet den befrier
fra Avtoritetstroens Forudsætninger, har dog netop
sin Grund i, at Naturideen hos ham, saameget den end er
Basis for Alt, staar som en fjærn og mystisk Magt og ikke
drages frem til konkret Gennemførelse i Virkeligheden.

Men det er naturligvis en ensidig Betragtningsmaade
at se Montaigne som filosofisk Etiker. Han kommer
kun til sin Ret, naar man ser ham i hans hele literær%
% --- p. 230 ---
\setcounter{footnote}{0}%
\opage{230}historiske Betydning, og denne beror paa, at han ikke
henvender sig til de Lærde og de ensomme Tænkere,
men at han i et saa let, klart og livligt Foredrag, som
endnu ikke var hørt i et nyere Sprog, rettede Talen til
Enhver, som vilde høre ham, om ikke af anden Grund,
saa dog for at underholde sig ved et saa behageligt Selskab
som Montaignes Essays er.
\signature{\textbf{Harald Høffding.}}

\end{document}
